\documentclass[10pt]{revtex4-2}
\usepackage{amsmath,amssymb,braket,newtxtext}
\usepackage[hidelinks]{hyperref}
\usepackage{color}
\usepackage{xcolor}
% \usepackage[margin=1.5cm,bottom=1.5cm]{geometry}
\newcommand{\shat}{\hat{\mathbf{s}}}
\newcommand{\question}[1]{\par \noindent {\bf Question.\quad}{#1}\\}
\newcommand{\response}[1]{\par \noindent {\bf Response.\quad}{\color{blue}{#1}}\\[1em]}

\begin{document}

\title{RESPONSE TO REVIEWER REPORTS FOR PHD THESIS}
\author{Abhirup Mukherjee (18IP014)}
\maketitle
\vspace{-20pt}
\hrulefill
\vspace{10pt}

\section*{Response to Prof. Pinaki Majumdar}

\question{One qualitative question: while the 1 + 4 cluster, for example, and the associated translation
operators preserve the symmetry of the 2D lattice (modulo any long range order), does
the process adopted generate the exact Greens functions of the lattice model?}

\response{
	In principle, the process lays out a hierarchy of contributions to the Greens function that would capture the lattice Greens function exactly. In practise, we have restricted the calculations to the first few terms; this is supported by the local nature of the impurity correlation that helps the higher terms drop off. The method of accounting for the additional terms are shown in eqs. 4.24 and 4.25.
}

\question{Maybe one could have had simpler notation for the electronic operators, I struggled to figure out what \(c_{Z,\sigma,f}\) might be.}
\response{We have now better explained the notation within the thesis.
}

\question{The standard model for heavy fermions is the Anderson lattice. If eqn3.2 is the lattice model here, I guess the Anderson lattice corresponds to Ud = 0, W = 0 and J = 0. Yet the phase diagrams are plotted in terms of these parameters. How do we connect to Anderson lattice results?
}

\response{We first note that the phase diagrams are plotted as a function of \(J_\perp\) (the inter-orbital Kondo coupling) and \(W_f\) (the \(f-\)layer bath correlation). The bath correlation for the conduction layer is set to zero (to make it only weakly interacting, compared to the \(f-\)layer). The phase diagrams are then plotted as a function of \(J_\perp\) and \(W_f\), to show how the behaviour changes upon tuning (i) the inter-layer connection, and (ii) the correlation of the \(f-\)layer. The Kondo lattice problem resides at large negative \(W_f\), because that's when the electrons of the \(f-\)layer form local moments. The presence of the coupling \(J_\perp\) implies that we are dealing with an extended periodic Anderson mode where the inter-orbital connectivity arises not just from one particle hybridisation but also from spin-exchange processes.
}

\question{There is detailed spectral work by D. E. Logan and coworkers on DMFT of the
Anderson lattice using a local moment approximation. Do the spectral results here
have any correspondence with them?}

\response{While Logan and coworkers have a series of works on the periodic Anderson model, we focus on~\cite{Smith2003}, where they deal with the particle-hole symmetric case that is pertinent to this thesis. The results obtained here agree qualitatively with those obtained in~\cite{Smith2003}, in terms of the opening of the gap and the movement of the Hubbard satellite peaks; more work is necessary in order to compare quantitative claims such as scaling of certain quantities as a function of frequency.
}

\section*{Response to Prof. Subhendra D. Mahanti}
\question{What is the physical difference between “Insulator with spin mixing between impurity spin and bath spin” and “Mott Insulator”. Explain}

\question{What is the physics of charge isospin Kondo coupling 𝐾 between the impurity and the bath: needs to be explained better.}

\response{
We first note that the effects of a charge isospin Kondo coupling hase not been explored within the thesis, and have only been mentioned in passing.

While the usual Kondo effect describes the screening of an impurity magnetic moment by the local spin density of the conduction bath, one can also envisage a scenario where it is the charge degrees of freedom of the impurity that is screened by the bath~\cite{taraphder1991}. This can happen, for example, if the local correlation on the impurity site is negative. That makes the \(n=0\) and \(n=2\) configurations (unoccupied and doubly occupied) lower in energy. One can then define an {\it isospin} from this doublet of states, as well as a similar isospin for the doublet of states in the bath. Isospin-flip fluctuations will then involve transitions where an unoccupied state on the impurity scatters to a double occupied state while the bath simultaneously scatters in the opposite manner. It is worth noting that during the Schrieffer-Wolff transformation from the Anderson impurity model to the Kondo model, such charge Kondo terms are also generated but they reside in the high-energy part of the effective theory~\cite{phillips2012advanced}.
}

\question{I have a problem with chapter 4 section 4.2.2, Periodization of Auxiliary Model into Extended
Hubbard Model using the lattice translation operator on the impurity cluster model. The operator \(c_{r_d,\sigma}\) is a destruction operator for an impurity electron located at the lattice site \(r_d\) (originally site
0 of the impurity model). Also, in the impurity model there is a destruction operator associated
with bath electrons associated with the same site, \(c_{0,\sigma}\). There are N such bath destruction
operators for N lattice sites. The number of fermion destruction operators is therefore N+1.
When one translates the impurity over the N lattice sites there are 2N fermion destruction
operators, not N as given in equations 2.6 and 2.7. What am I missing. I would have expected a
2-band extended Hubbard model, one for the “impurity” and the other for the bath.}

\response{If the total number of lattice sites is \(N\), then the number of bath operators is \(N-1\) because the impurity site is "embedded" in the lattice and forms the \(N^\text{th}\) operator. This is necessary in order for the impurity couplings to be sensitive to the lattice details. Applying translation operators does not double the number of operators because the same operators transform into one another: \(c(r) \to c(r + a), c(r + a) \to c(r + 2a) \ldots c(r+Na) \to c(r)\).}

\question{All the discussions are associated with an extended 2D Hubbard model. Are you assuming
that the solution of this Hamiltonian is simply the periodic translation of the solution of the
impurity embedded in a 2D lattice. Is this an ansatz or are there some justifications for this
approach?}

\response{Indeed, there's a set of arguments that relate the impurity model to the extended Hubbard model on the 2D square lattice. First, we show that the lattice model can be obtained by repeated translation of the impurity model (eq. 2.11). This motivates a form of the eigenstates of the lattice model in terms of those of the impurity model (eq. 2.20), via a many-body Blochs's theorem. Combining these two properties, we show that the lattice model Greens function can be expressed in terms of auxiliary mode quantities (eq. 4.23). This final relation is what allows us to reconstruct the low-energy physics of the lattice model.}

\question{The novel feature is the condition of sufficiently large attractive W. It is not clear to me how such an attractive interaction appears in a simple 2D Hubbard model. But the most interesting observation is the appearance of a third fixed point, pseudo gap fixed point in a finite range of W.}

\response{We feel that the attractive bath interaction \(W \) is best not understood as a microscopic coupling, but rather as an emergent effective coupling. It captures the tendency of the conduction electrons to favour unoccupied or double occupied states (relative to the singly-occupied states). Perhaps in a more illuminating sense, just as the Coulomb repulsion \(U\) represents the tendency that competes with the hopping processes in the simple Hubbard model, the bath correlation \(W\) represents the tendency that competes with the local Fermi liquid excitations that propagate out of the impurity site. In this sense, it is a proxy for the hubbard U in the neighbourhood of the impurity site; both compete against delocalisation processes.}

\question{{\color{red}Is the pseudo gap fixed point associated with insulating and spin mixing as discussed in the previous chapter.} Can you describe the spin density distribution and spin-spin correlation
function associated with this fixed point. What role does the lattice symmetry play in the
nature of this pseudo gap fixed point. What will happen if the 2D lattice is triangular or
hexagonal?}

\response{The unitary RG, as a matter of principle, does not allow for spontaneous symmetry breaking (SSB) because unitary transformations preserve all symmetries. One way of searching for SSB phases within such a formalism is to put in an explicit symmetry breaking term in the Hamiltonian and then study whether it is relevant under RG flow. We haven't carried out such an analysis within this thesis, and as a result, we don't find any such phase. Accordingly, in the absence of any magnetic order, there is no spin density in the system in any phase. The spin-spin correlation, however, shows long-ranged behaviour in the pseudogap phase (Fig. 5.13).

\par With regards to the role of lattice symmetry, we haven't checked whether our conclusions hold in other dimensions or other lattices. It sounds reasonable that changing the lattice symmetries will affect many of the results. However, a recently growing body of work~\cite{Huang2022,lanave2025,bai2026localcurrentalgebrahk} suggests that the result that the Hatsugai-Kohmoto model is the parent metal of the Mott insulator can be grounded on more robust ideas such as the breaking of generalised symmetries and anomalies of the Luttinger-Ward functional. These indicate that our results are likely to extend beyond just the 2D square lattice.
}

\question{I had some problem with the strong conclusion “Thus, the drastic change in nature of the
real-space excitations- from locally well-defined Landau quasiparticles of the FL to the
increasingly nonlocal excitations of the NFL Fermi arcs in the PG phase-appears to be
dictated by a topological principle and, therefore, robust under renormalization. This
signals the NFL Fermi arcs of the PG as an emergent phase of strongly interacting
quantum matter- which we dub the Mott metal- that is the parent metal of the MI.”
What is the role of lattice dimensionality (D=2 vs 3) and lattice geometry (square vs triangular or Kagome’) on this conclusion?
}

\response{
	As mentioned in the previous response, very strictly speaking our results apply only to the case of the 2D square lattice. However, a recently growing body of work~\cite{Huang2022,lanave2025,bai2026localcurrentalgebrahk} suggests that the result that the Hatsugai-Kohmoto model is the parent metal of the Mott insulator can be grounded on more robust ideas such as the breaking of generalised symmetries and anomalies of the Luttinger-Ward functional. These indicate that our results are likely to extend beyond just the 2D square lattice.
}

\question{What happens for large \(W_f\) and one increases \(J_\perp\)?}

\response{
	For large \(W_f\), the \(f-\)layer becomes a Mott insulator. The pseudogap phase shrinks to a smaller region due to the lack of non-Fermi liquid behaviour in the \(f-\)layer. Increasing \(J_\perp\) brings about a Kondo insulator.
}

\question{He thinks the appearance of long-range correlation at the transition from Fermi liquid to a different phase might be relevant to the observation of strange metallic behavior in heavy-fermion systems at high temperatures above the quantum critical point. But all the calculations are at T=0K.}

\response{
	While it is true that our results are all at \(T=0\), it stands to reason that some of these phases will have finite-temperature effects. For example, the pseudogap non-Fermi liquid behaviour we see near the transition arises from the quantum critical fluctuations; these fluctuations can survive up to certain temperatures and have finite temperature counterparts.
}

\section{Additional changes}
We have made the following additional changes to the manuscript:
\begin{itemize}
	\item We have improved the inset figure of panel (a) in Fig. 4 for the impurity magnetisation (computed earlier for a smaller system size (49 x 49)) to that for a larger system size (77 x 77, same as the rest of the manuscript). The updated plot shows more clearly the dramatic rise of the impurity magnetisation towards the end of the pseudogap phase, and upon approaching the Mott critical point. This highlights the near complete breakdown of Kondo screening at Mott criticality as observed by us.
	\item We have added references to two works~(\cite{si_kotliar_1993,kotliarsi_1993}) from the past that adopt an impurity model-based approach towards heavy fermions that are philosophically similar to ours.
	\item We have also fixed a typographical error below eq. (36) regarding the definition of the operator \(r_{q\sigma}\).
\end{itemize}

\bibliography{tilingProject.bib}
\bibliographystyle{apsrev4-2}
\end{document}
